% Where a sample comes from and what happens to it. Columns on a 2.2 cm pitch.
\begin{tikzpicture}[font=\sffamily\small, x=1cm, y=1cm]

% ---------------- the sensor, on its own clock
\node[hw, text width=32mm] (mems) at (-6.0,3.0)
  {inertial unit\\{\scriptsize rate and acceleration, on the shield}};
\node[hw, text width=32mm] (fifo) at (-6.0,0.8)
  {its buffer\\{\scriptsize N samples deep, its own oscillator}};
\node[hw, text width=32mm] (int)  at (-6.0,-1.4)
  {data-ready pin\\{\scriptsize raised when the watermark is reached}};
\begin{scope}[on background layer]
\node[layer, fit=(mems)(fifo)(int), inner sep=6pt,
      label={[layerlabel, anchor=south west]north west:the sensor's own time}] {};
\end{scope}

% ---------------- the node
\node[fw, text width=34mm] (cap)  at (0.2,-1.4)
  {timer capture\\{\scriptsize one timestamp, taken at the pin, in hardware}};
\node[fw, text width=34mm] (task) at (0.2,0.8)
  {reader task\\{\scriptsize one burst read per batch, never in the period handler}};
\node[fw, text width=34mm] (stamp) at (0.2,3.0)
  {reconstruct and rotate\\{\scriptsize a timestamp per sample, then into the joint's axes}};
\begin{scope}[on background layer]
\node[layer, fit=(cap)(task)(stamp), inner sep=6pt,
      label={[layerlabel, anchor=south west]north west:the node's time}] {};
\end{scope}

% ---------------- what it feeds
\node[user, text width=30mm] (ring)  at (6.2,3.0) {sample ring\\{\scriptsize each sample with its own timestamp and its age}};
\node[user, text width=30mm] (loop)  at (6.2,0.8) {the period handler\\{\scriptsize takes the nearest sample, performs no transaction}};
\node[user, text width=30mm] (cross) at (6.2,-1.4) {cross-check\\{\scriptsize against the encoder from chapter 5}};

% ---------------- what an inertial unit cannot give
\node[absentblk, text width=32mm] (angle) at (0.2,-3.7)
  {the joint angle\\{\scriptsize one unit on one link cannot measure it at all. It
   measures the link's motion in space, and gravity}};

\draw[flow] (mems) -- (fifo);
\draw[flow] (fifo) -- (int);
\draw[flow] (int)  -- node[lbl, above]{pin} (cap);
\draw[flow] (fifo.east) -- ++(1.1,0) |- node[lbl, above, pos=0.3]{burst read} (task.west);
\draw[flow] (cap)  -- (task);
\draw[flow] (task) -- (stamp);
\draw[flow] (stamp) -- (ring);
\draw[flow] (ring) -- (loop);
\draw[flow] (loop) -- (cross);
\draw[hiz, ->] (angle.east) -- ++(1.4,0) |- node[lbl, right, pos=0.7]{not this} (cross.south);

\node[umlnote, text width=52mm, anchor=north west] at (-9.6,-3.5)
  {The hardware gives exactly one timestamp per batch: the capture at the
   data-ready pin. Everything else in this figure is arithmetic on that one
   number, with an error bounded by how well the nominal sample interval matches
   the interval the sensor actually achieves, which is why the chapter measures
   that interval instead of trusting it.};
\end{tikzpicture}
